% !TeX program = xelatex
\documentclass[11pt,a4paper,fontset=fandol]{ctexart}
\usepackage{guide}

\title{我的研究笔记}
\subtitle{记录问题、方法与下一步行动}
\author{你的名字}
\institute{课题组 / 单位}
\date{\today}
\guideheader{研究笔记}{项目名称}
% \guidetheme{royalpurple}

\begin{document}
\maketitle
% 长文可启用目录：
% \tableofcontents
% \clearpage

\iconsection{bullseye}{问题与目标}
在这里写正文。用 \important{红色粗体} 标注重点，
用 \themecolor{主题色} 突出关键词，或用 \highlight{浅黄色底纹} 标记短语。

\emoji{lightbulb}~一个值得继续探索的想法。

\begin{tipbox}[核心想法]
先说明要解决的问题，再写自己的方法和预期结果。
\end{tipbox}

\iconsection{tasks}{当前进展}
\begin{itemize}
  \gooditem 已完成：写下本周取得的进展。
  \warningitem 需要关注：写下仍需解决的问题。
  \arrowitem 下一步：写下明确、可执行的行动。
\end{itemize}

\begin{infobox}[记录方式]
每次更新保留关键结论、对应证据和下一步行动。
\end{infobox}
\end{document}
